\question
Кот Барсик учит хозяйку своему языку. В нём три звука: «мр», «мяу» и «фыр». Каждому слову кошачьего языка Барсик сопоставляет код -- последовательность звуков; во фразе коды произносятся подряд, без пауз. Барсик хочет, чтобы любую фразу можно было расшифровать однозначно. Длины кодов уже придуманных слов: $1$, $2$, $2$ (сами коды ещё не выбраны). Какое наибольшее число новых слов с кодами длины не больше $3$ может добавить Барсик? Ответ обоснуйте.

---------------

Автор -- Дарья Гаркуша